% BEGIN VOL07-EXPANSION VII31
\section{Supplementary worked examples}
\label{sec:vii31-pedagogy}

\begin{example}[Hyperboloid point to disk]\label{ex:vii31-ped-01}
The hyperboloid point \((0,0,1)\) maps under \(\Phi\) to \(0\in\mathbb D\).  More generally, points with large \(x_3\) project close to the unit circle, making the ideal boundary visible as an asymptotic limit.
\end{example}

\begin{example}[Cayley normalization]\label{ex:vii31-ped-02}
The Cayley transform \(C(z)=i(1+z)/(1-z)\) sends \(0\) to \(i\), \(-1\) to \(0\), and the boundary point \(1\) to \(\infty\).  One formula therefore aligns the disk's ideal circle with the real-line boundary of the half-plane.
\end{example}

\begin{example}[Geodesic pictures]\label{ex:vii31-ped-03}
The same hyperbolic geodesic appears as a plane section in the hyperboloid model, a circle orthogonal to the boundary in the Poincare disk, and a chord in the Klein disk.  Model choice changes coordinate appearance, not intrinsic geometry.
\end{example}

\section{Graded supplementary exercises}
The following exercises add a deliberate progression from concrete calculation to structural proof, hypothesis testing, application, and synthesis.

\subsection*{Standard computations and constructions}

\begin{exercise}[Base point]\label{exr:vii31-ped-01}
Where does \((0,0,1)\) on the hyperboloid map in the disk model?
\end{exercise}

\begin{hint}
Use \(\Phi=(x_1+ix_2)/(x_3+1)\).
\end{hint}

\begin{solution}
It maps to \(0\).
\end{solution}

\begin{exercise}[Inverse base point]\label{exr:vii31-ped-02}
Compute \(\Phi^{-1}(0)\).
\end{exercise}

\begin{hint}
Insert \(z=0\).
\end{hint}

\begin{solution}
\((0,0,1)\).
\end{solution}

\begin{exercise}[Cayley origin]\label{exr:vii31-ped-03}
Compute \(C(0)\).
\end{exercise}

\begin{hint}
Substitute directly.
\end{hint}

\begin{solution}
\(i\).
\end{solution}

\begin{exercise}[Cayley inverse]\label{exr:vii31-ped-04}
Compute \(C^{-1}(i)\).
\end{exercise}

\begin{hint}
Use \((w-i)/(w+i)\).
\end{hint}

\begin{solution}
\(0\).
\end{solution}

\begin{exercise}[Ideal boundary]\label{exr:vii31-ped-05}
What is the ideal boundary of the upper half-plane model?
\end{exercise}

\begin{hint}
Add the point at infinity to the real line.
\end{hint}

\begin{solution}
\(\mathbb R\cup\{\infty\}\).
\end{solution}

\subsection*{Proofs}

\begin{exercise}[Tangent positivity]\label{exr:vii31-ped-06}
Why is the Lorentz form positive definite on the tangent space of the hyperboloid?
\end{exercise}

\begin{hint}
The position vector is timelike.
\end{hint}

\begin{solution}
The orthogonal complement of a timelike vector in signature \((2,1)\) is spacelike and positive definite.
\end{solution}

\begin{exercise}[Model map]\label{exr:vii31-ped-07}
Verify \(\Phi^{-1}(z)\) has Lorentz norm \(-1\).
\end{exercise}

\begin{hint}
Substitute the explicit inverse coordinates.
\end{hint}

\begin{solution}
The numerator simplifies to \(-\left(1-|z|^2\right)^2\), which divided by the same squared denominator gives \(-1\).
\end{solution}

\begin{exercise}[Disk range]\label{exr:vii31-ped-08}
Why does \(\Phi\) land inside \(|z|<1\)?
\end{exercise}

\begin{hint}
Use \(x_3^2-x_1^2-x_2^2=1\).
\end{hint}

\begin{solution}
One obtains \(|\Phi|^2=(x_3-1)/(x_3+1)<1\) since \(x_3>0\).
\end{solution}

\begin{exercise}[Model equivalence]\label{exr:vii31-ped-09}
What does it mean for two hyperbolic models to be equivalent?
\end{exercise}

\begin{hint}
The identifying map must preserve the hyperbolic metric.
\end{hint}

\begin{solution}
They are related by an isometry, so all intrinsic lengths, angles, geodesics, and curvatures agree after translation through the map.
\end{solution}

\subsection*{Counterexamples and hypothesis tests}

\begin{exercise}[Boundary belongs to space]\label{exr:vii31-ped-10}
Test the claim for the disk model.
\end{exercise}

\begin{hint}
The model is the open unit disk.
\end{hint}

\begin{solution}
False.  The unit circle is an ideal boundary, not part of \(\mathbb H^2\).
\end{solution}

\begin{exercise}[Klein angles]\label{exr:vii31-ped-11}
Test the claim: every standard disk model preserves Euclidean angles.
\end{exercise}

\begin{hint}
Distinguish Poincare and Klein models.
\end{hint}

\begin{solution}
False.  The Poincare disk is conformal; the Klein disk is not.
\end{solution}

\begin{exercise}[Model geodesics look same]\label{exr:vii31-ped-12}
Test the claim: a geodesic has the same Euclidean shape in every model.
\end{exercise}

\begin{hint}
Compare Poincare circles and Klein chords.
\end{hint}

\begin{solution}
False.  Model transformations change coordinate shapes while preserving intrinsic geodesicity.
\end{solution}

\subsection*{Applications and investigations}

\begin{exercise}[Visualization]\label{exr:vii31-ped-13}
Why is the Klein model useful for visualizing geodesic incidence?
\end{exercise}

\begin{hint}
Its geodesics are Euclidean chords.
\end{hint}

\begin{solution}
Straight chord geometry makes intersections and convexity visually simple, even though angles are distorted.
\end{solution}

\begin{exercise}[Angle work]\label{exr:vii31-ped-14}
Why is the Poincare model preferable for angle calculations?
\end{exercise}

\begin{hint}
It is conformal.
\end{hint}

\begin{solution}
Hyperbolic and Euclidean angles coincide at intersection points.
\end{solution}

\subsection*{Challenge problems}

\begin{exercise}[Boundary limit]\label{exr:vii31-ped-15}
What happens to the disk coordinate of a hyperboloid point as \(x_3\to\infty\) along a fixed asymptotic direction?
\end{exercise}

\begin{hint}
Use \(|\Phi|^2=(x_3-1)/(x_3+1)\).
\end{hint}

\begin{solution}
Its modulus tends to \(1\), approaching the ideal boundary.
\end{solution}

\begin{exercise}[Cayley endpoints]\label{exr:vii31-ped-16}
Which disk boundary point maps to \(\infty\) under \(C\)?
\end{exercise}

\begin{hint}
Find where the denominator vanishes.
\end{hint}

\begin{solution}
\(z=1\).
\end{solution}

% END VOL07-EXPANSION VII31
